\section*{الفصل \ref{ch:g7:identifying-substances} --- التعرف على المواد}

\begin{solution}{exo:g7:identifying-substances:1}
اختبار ماء الجير: يُمرّر الغاز فقاعاتٍ في ماء الجير الرائق، فيتعكّر إذا كان
الغاز ثنائي أكسيد الكربون.
\end{solution}

\begin{solution}{exo:g7:identifying-substances:2}
المسحوق هو كبريتات النحاس اللامائية، والسائل يحوي ماءً.
\end{solution}

\begin{solution}{exo:g7:identifying-substances:3}
الأكسجين.
\end{solution}

\begin{solution}{exo:g7:identifying-substances:4}
رصاص القلم لا \omterm{def:g2:mixing-and-dissolving:dissolve}{يذوب} في \omterm{def:g6:solutions-and-solubility:solute}{المذيب}. أما خط الحبر فسيصعد هو نفسه وينفصل،
فيفسد \omterm{def:g7:identifying-substances:chromatography}{الكروماتوغرام}.
\end{solution}

\begin{solution}{exo:g7:identifying-substances:5}
لو كانت البقعة تحت \omterm{def:g6:solutions-and-solubility:solute}{المذيب} لذابت الصبغات في \omterm{def:g6:solutions-and-solubility:solute}{مذيب} الكأس بدلًا من أن
تصعد في الورقة.
\end{solution}

\begin{solution}{exo:g7:identifying-substances:6}
لا، إنه \omterm{def:g6:pure-substances-and-mixtures:mixture}{خليط} من صبغتين: على الأرجح الصبغة الصفراء والصبغة الزرقاء
المرجعيتان.
\end{solution}

\begin{solution}{exo:g7:identifying-substances:7}
اجمع قليلًا من الغاز في أنبوب اختبار وأمسك عودًا مشتعلًا عند فوهته: فالفرقعة
الحادة تدل على الهيدروجين.
\end{solution}

\begin{solution}{exo:g7:identifying-substances:8}
\omterm{def:g4:air-a-mixture-of-gases:air}{الهواء} المزفور يحوي من ثنائي أكسيد الكربون أكثر بكثير مما يحويه \omterm{def:g4:air-a-mixture-of-gases:air}{الهواء} النقي.
\end{solution}

\begin{solution}{exo:g7:identifying-substances:9}
يبيّن الاختبار وجود الماء، لا أنه وحده: فالماء المالح أو العصير أو أي \omterm{def:g6:pure-substances-and-mixtures:mixture}{خليط}
مائي يُزرّق المسحوق أيضًا.
\end{solution}

\begin{solution}{exo:g7:identifying-substances:10}
ثلاث صبغات (ثلاث بقع). والصبغة الصفراء صعدت أعلى من غيرها.
\end{solution}

\begin{solution}{exo:g7:identifying-substances:11}
مثلًا: عود متوهج في كل مخبار --- فالمخبار الذي يعيد إشعاله فيه الأكسجين.
وفي المخبارين الآخرين، عود مشتعل عند الفوهة --- فالفرقعة الحادة تدل على
الهيدروجين. والمخبار الأخير فيه ثنائي أكسيد الكربون (ويمكن التحقق منه
بماء الجير). فاختباران أو ثلاثة تكفي.
\end{solution}

\begin{solution}{exo:g7:identifying-substances:12}
قد لا يحوي الغاز ثنائي أكسيد الكربون، أو يحوي منه قليلًا جدًا لا يظهر في
تلك المدة. واختبار شاهد \omterm{def:g4:air-a-mixture-of-gases:air}{بهواء} الزفير، المعروف أنه يحوي ثنائي أكسيد
الكربون، يتحقق من أن ماء الجير صالح للعمل.
\end{solution}

\begin{solution}{pb:g7:identifying-substances:1}
\textbf{1.} لكي تكون الشروط (الورق، \omterm{def:g6:solutions-and-solubility:solute}{والمذيب}، والمدة) واحدة للأحبار
الأربعة: فعندئذ فقط يمكن مقارنة ارتفاعات البقع.

\textbf{2.} \omterm{def:g6:pure-substances-and-mixtures:mixture}{خليط}: فحبرها يعطي ثلاث بقع منفصلة، أي ثلاث صبغات.

\textbf{3.} صبغتان. لا: فحبر الورقة فيه ثلاث صبغات، وليس في القلم 1 لا
الصبغة الحمراء ولا الصبغة الصفراء.

\textbf{4.} بقعه الثلاث ليست في ارتفاعات بقع الورقة نفسها: فهي ثلاث
صبغات مختلفة.

\textbf{5.} القلم 2: فبقعه الثلاث في ارتفاعات بقع الورقة الثلاث
تمامًا.

\textbf{6.} الأكسجين.

\textbf{7.} الهيدروجين.

\textbf{8.} ثنائي أكسيد الكربون.

\textbf{9.} اختبار شاهد (شاهد سلبي). وهو يبيّن أن ماء الجير لا يتعكّر
من تلقاء نفسه، ولا \omterm{def:g4:air-a-mixture-of-gases:air}{بالهواء} العادي في تلك المدة: فتعكّره بالغاز Z آتٍ حقًا
من ثنائي أكسيد الكربون.

\textbf{10.} «الماء --- كبريتات النحاس اللامائية --- يصير الأبيض أزرق»:
السائل يحوي ماءً.

\textbf{11.} لا. فالاختبار الإيجابي يبيّن أن الماء موجود، لا أنه لا يوجد
شيء غيره: فقد يكون السائل \omterm{def:g2:mixing-and-dissolving:dissolve}{محلولًا} أو خليطًا مائيًا آخر.

\textbf{12.} حبر الورقة يحوي 3 صبغات، والقلم 2 هو الذي كتب
الورقة.
\end{solution}
